%\title{University of Bristol Thesis Template}
\RequirePackage[l2tabu]{nag}		% Warns for incorrect (obsolete) LaTeX usage
%
%
% File: memoirthesis.tex
% Author: Victor Brena
% Description: Contains the thesis template using memoir class,
% which is mainly based on book class but permits better control of 
% chapter styles for example. This template is an adaptation and 
% modification of Oscar's.
% 
% Memoir is a flexible class for typesetting poetry, fiction, 
% non-fiction and mathematical works as books, reports, articles or
% manuscripts. CTAN repository is found at:
% http://www.ctan.org/tex-archive/macros/latex/contrib/memoir/
%
%
% UoB guidelines for thesis presentation were found at:
% http://www.bris.ac.uk/esu/pg/pgrcop11-12topic.pdf#page=49
%
% UoB guidelines:
%
% The dissertation must be printed on A4 white paper. Paper up to A3 may be used
% for maps, plans, diagrams and illustrative material. Pages (apart from the
% preliminary pages) should normally be double-sided.
%
% Memoir class loads useful packages by default (see manual).
\documentclass[a4paper,11pt,leqno,openbib]{memoir} %add 'draft' to turn draft option on (see below)
\setlength{\parindent}{0em}
\setlength{\parskip}{1em}
%
%
% Adding metadata:
\usepackage{datetime}
\usepackage{ifpdf}
\ifpdf
\pdfinfo{
   /Author (Eddy Bertoluzzo)
   /Title (Master Thesis)
   /Keywords (One; Two;Three)
   /CreationDate (D:\pdfdate)
}
\fi
% When draft option is on. 
\ifdraftdoc 
	\usepackage{draftwatermark}				%Sets watermarks up.
	\SetWatermarkScale{0.3}
	\SetWatermarkText{\bf Draft: \today}
\fi
%
% Declare figure/table as a subfloat.
\newsubfloat{figure}
\newsubfloat{table}
% Better page layout for A4 paper, see memoir manual.
\settrimmedsize{297mm}{210mm}{*}
\setlength{\trimtop}{0pt} 
\setlength{\trimedge}{\stockwidth} 
\addtolength{\trimedge}{-\paperwidth} 
\settypeblocksize{634pt}{448.13pt}{*} 
\setulmargins{4cm}{*}{*} 
\setlrmargins{*}{*}{1.5} 
\setmarginnotes{17pt}{51pt}{\onelineskip} 
\setheadfoot{\onelineskip}{2\onelineskip} 
\setheaderspaces{*}{2\onelineskip}{*} 
\checkandfixthelayout
%
\frenchspacing
% Font with math support: New Century Schoolbook
\usepackage{fouriernc}
\usepackage[T1]{fontenc}
%
% UoB guidelines:
%
% Text should be in double or 1.5 line spacing, and font size should be
% chosen to ensure clarity and legibility for the main text and for any
% quotations and footnotes. Margins should allow for eventual hard binding.
%
% Note: This is automatically set by memoir class. Nevertheless \OnehalfSpacing 
% enables double spacing but leaves single spaced for captions for instance. 
\OnehalfSpacing 
%
% Sets numbering division level
\setsecnumdepth{subsection} 
\maxsecnumdepth{subsubsection}
%
% Chapter style (taken and slightly modified from Lars Madsen Memoir Chapter 
% Styles document
\usepackage{calc,soul,fourier}
\makeatletter 
\newlength\dlf@normtxtw 
\setlength\dlf@normtxtw{\textwidth} 
\newsavebox{\feline@chapter} 
\newcommand\feline@chapter@marker[1][4cm]{%
	\sbox\feline@chapter{% 
		\resizebox{!}{#1}{\fboxsep=1pt%
			\colorbox{gray}{\color{white}\thechapter}% 
		}}%
		\rotatebox{90}{% 
			\resizebox{%
				\heightof{\usebox{\feline@chapter}}+\depthof{\usebox{\feline@chapter}}}% 
			{!}{\scshape\so\@chapapp}}\quad%
		\raisebox{\depthof{\usebox{\feline@chapter}}}{\usebox{\feline@chapter}}%
} 
\newcommand\feline@chm[1][4cm]{%
	\sbox\feline@chapter{\feline@chapter@marker[#1]}% 
	\makebox[0pt][c]{% aka \rlap
		\makebox[1cm][r]{\usebox\feline@chapter}%
	}}
\makechapterstyle{daleifmodif}{
	\renewcommand\chapnamefont{\normalfont\Large\scshape\raggedleft\so} 
	\renewcommand\chaptitlefont{\normalfont\Large\bfseries\scshape} 
	\renewcommand\chapternamenum{} \renewcommand\printchaptername{} 
	\renewcommand\printchapternum{\null\hfill\feline@chm[2.5cm]\par} 
	\renewcommand\afterchapternum{\par\vskip\midchapskip} 
	\renewcommand\printchaptertitle[1]{\color{gray}\chaptitlefont\raggedleft ##1\par}
} 
\makeatother 
\chapterstyle{daleifmodif}
%
% UoB guidelines:
%
% The pages should be numbered consecutively at the bottom centre of the
% page.
\makepagestyle{myvf} 
\makeoddfoot{myvf}{}{\thepage}{} 
\makeevenfoot{myvf}{}{\thepage}{} 
\makeheadrule{myvf}{\textwidth}{\normalrulethickness} 
\makeevenhead{myvf}{\small\textsc{\leftmark}}{}{} 
\makeoddhead{myvf}{}{}{\small\textsc{\rightmark}}
\pagestyle{myvf}
%
% Oscar's command (it works):
% Fills blank pages until next odd-numbered page. Used to emulate single-sided
% frontmatter. This will work for title, abstract and declaration. Though the
% contents sections will each start on an odd-numbered page they will
% spill over onto the even-numbered pages if extending beyond one page
% (hopefully, this is ok).
\newcommand{\clearemptydoublepage}{\newpage{\thispagestyle{empty}\cleardoublepage}}
%
%
% Creates indexes for Table of Contents, List of Figures, List of Tables and Index
\makeindex
% \printglossaries below creates a list of abbreviations. \gls and related
% commands are then used throughout the text, so that latex can automatically
% keep track of which abbreviations have already been defined in the text.
%
% The import command enables each chapter tex file to use relative paths when
% accessing supplementary files. For example, to include
% chapters/brewing/images/figure1.png from chapters/brewing/brewing.tex we can
% use
% \includegraphics{images/figure1}
% instead of
% \includegraphics{chapters/brewing/images/figure1}
\usepackage{import}

% Add other packages needed for chapters here. For example:
\usepackage{lipsum}					%Needed to create dummy text
\usepackage{amsfonts} 					%Calls Amer. Math. Soc. (AMS) fonts
\usepackage[centertags]{amsmath}			%Writes maths centred down
\usepackage{stmaryrd}					%New AMS symbols
\usepackage{amssymb}					%Calls AMS symbols
\usepackage{amsthm}					%Calls AMS theorem environment
\usepackage{newlfont}					%Helpful package for fonts and symbols
\usepackage{layouts}					%Layout diagrams
\usepackage{graphicx}					%Calls figure environment
\usepackage{longtable,rotating}			%Long tab environments including rotation. 
\usepackage[applemac]{inputenc}			%Needed to encode non-english characters 
									%directly for mac
\usepackage{colortbl}					%Makes coloured tables
\usepackage{wasysym}					%More math symbols
\usepackage{mathrsfs}					%Even more math symbols
\usepackage{float}						%Helps to place figures, tables, etc. 
\usepackage{verbatim}					%Permits pre-formated text insertion
\usepackage{upgreek }					%Calls other kind of greek alphabet
\usepackage{latexsym}					%Extra symbols
\usepackage[square,numbers,
		     sort&compress]{natbib}		%Calls bibliography commands 
\usepackage{url}						%Supports url commands
\usepackage{etex}						%eTeXÕs extended support for counters
\usepackage{fixltx2e}					%Eliminates some in felicities of the 
									%original LaTeX kernel
\usepackage[spanish,english]{babel}		%For languages characters and hyphenation
\usepackage{color}                    				%Creates coloured text and background
\usepackage[colorlinks=true,
		     allcolors=black]{hyperref}              %Creates hyperlinks in cross references
\usepackage{memhfixc}					%Must be used on memoir document 
									%class after hyperref
\usepackage{enumerate}					%For enumeration counter
\usepackage{footnote}					%For footnotes
\usepackage{microtype}					%Makes pdf look better.
\usepackage{rotfloat}					%For rotating and float environments as tables, 
									%figures, etc. 
\usepackage{alltt}						%LaTeX commands are not disabled in 
									%verbatim-like environment
\usepackage[version=0.96]{pgf}			%PGF/TikZ is a tandem of languages for producing vector graphics from a 
\usepackage{tikz}						%geometric/algebraic description.
\usetikzlibrary{arrows,shapes,snakes,
		       automata,backgrounds,
		       petri,topaths}				%To use diverse features from tikz		
%		

% Todo package
\usepackage[colorinlistoftodos,prependcaption,textsize=tiny]{todonotes}
%

%Minted package/ Code formatting
\usepackage{xcolor}
\usepackage{minted}
%\usemintedstyle{colorful}	

\definecolor{mintedbackground}{rgb}{0.95,0.95,0.95}
\newcommand{\code}[1]{\colorbox{mintedbackground}{\texttt{#1}}}

\newmintedfile[haskellcode]{haskell}{
bgcolor=mintedbackground,
linenos=true,
numberblanklines=true,
numbersep=12pt,
numbersep=5pt,
gobble=0,
frame=leftline,
framerule=0.4pt,
framesep=2mm,
funcnamehighlighting=true,
tabsize=4,
obeytabs=false,
mathescape=false
samepage=false, %with this setting you can force the list to appear on the same page
showspaces=false,
showtabs =false,
texcl=false
}
				

%Reduce widows  (the last line of a paragraph at the start of a page) and orphans 
% (the first line of paragraph at the end of a page)
\widowpenalty=1000
\clubpenalty=1000
%
% New command definitions for my thesis
%
\newcommand{\keywords}[1]{\par\noindent{\small{\bf Keywords:} #1}} %Defines keywords small section
\newcommand{\parcial}[2]{\frac{\partial#1}{\partial#2}}                             %Defines a partial operator
\newcommand{\vectorr}[1]{\mathbf{#1}}                                                        %Defines a bold vector
\newcommand{\vecol}[2]{\left(                                                                         %Defines a column vector
	\begin{array}{c} 
		\displaystyle#1 \\
		\displaystyle#2
	\end{array}\right)}
\newcommand{\mados}[4]{\left(                                                                       %Defines a 2x2 matrix
	\begin{array}{cc}
		\displaystyle#1 &\displaystyle #2 \\
		\displaystyle#3 & \displaystyle#4
	\end{array}\right)}
\newcommand{\pgftextcircled}[1]{                                                                    %Defines encircled text
    \setbox0=\hbox{#1}%
    \dimen0\wd0%
    \divide\dimen0 by 2%
    \begin{tikzpicture}[baseline=(a.base)]%
        \useasboundingbox (-\the\dimen0,0pt) rectangle (\the\dimen0,1pt);
        \node[circle,draw,outer sep=0pt,inner sep=0.1ex] (a) {#1};
    \end{tikzpicture}
}
\newcommand{\range}[1]{\textnormal{range }#1}                                             %Defines range operator
\newcommand{\innerp}[2]{\left\langle#1,#2\right\rangle}                                 %Defines inner product
\newcommand{\prom}[1]{\left\langle#1\right\rangle}                                         %Defines average operator
\newcommand{\tra}[1]{\textnormal{tra} \: #1}                                                       %Defines trace operator
\newcommand{\sign}[1]{\textnormal{sign\,}#1}                                                   %Defines sign operator
\newcommand{\sech}[1]{\textnormal{sech} #1}                                                  %Defines sech
\newcommand{\diag}[1]{\textnormal{diag} #1}                                                    %Defines diag operator
\newcommand{\arcsech}[1]{\textnormal{arcsech} #1}                                       %Defines arcsech
\newcommand{\arctanh}[1]{\textnormal{arctanh} #1}                                         %Defines arctanh
%Change tombstone symbol
\newcommand{\blackged}{\hfill$\blacksquare$}
\newcommand{\whiteged}{\hfill$\square$}
\newcounter{proofcount}
\renewenvironment{proof}[1][\proofname.]{\par
 \ifnum \theproofcount>0 \pushQED{\whiteged} \else \pushQED{\blackged} \fi%
 \refstepcounter{proofcount}
 \normalfont 
 \trivlist
 \item[\hskip\labelsep
       \itshape
   {\bf\em #1}]\ignorespaces
}{%
 \addtocounter{proofcount}{-1}
 \popQED\endtrivlist
}
%
%
% New definition of square root:
% it renames \sqrt as \oldsqrt
\let\oldsqrt\sqrt
% it defines the new \sqrt in terms of the old one
\def\sqrt{\mathpalette\DHLhksqrt}
\def\DHLhksqrt#1#2{%
\setbox0=\hbox{$#1\oldsqrt{#2\,}$}\dimen0=\ht0
\advance\dimen0-0.2\ht0
\setbox2=\hbox{\vrule height\ht0 depth -\dimen0}%
{\box0\lower0.4pt\box2}}
%
% My caption style
\newcommand{\mycaption}[2][\@empty]{
	\captionnamefont{\scshape} 
	\changecaptionwidth
	\captionwidth{0.9\linewidth}
	\captiondelim{.\:} 
	\indentcaption{0.75cm}
	\captionstyle[\centering]{}
	\setlength{\belowcaptionskip}{10pt}
	\ifx \@empty#1 \caption{#2}\else \caption[#1]{#2}
}
%
% My subcaption style
\newcommand{\mysubcaption}[2][\@empty]{
	\subcaptionsize{\small}
	\hangsubcaption
	\subcaptionlabelfont{\rmfamily}
	\sidecapstyle{\raggedright}
	\setlength{\belowcaptionskip}{10pt}
	\ifx \@empty#1 \subcaption{#2}\else \subcaption[#1]{#2}
}
%
%An initial of the very first character of the content
\usepackage{lettrine}
\newcommand{\initial}[1]{%
	\lettrine[lines=3,lhang=0.33,nindent=0em]{
		\color{gray}
     		{\textsc{#1}}}{}}
%
% Theorem styles used in my thesis
%
\theoremstyle{plain}
\newtheorem{theo}{Theorem}[chapter]
\theoremstyle{plain}
\newtheorem{prop}{Proposition}[chapter]
\theoremstyle{plain}
\theoremstyle{definition}
\newtheorem{dfn}{Definition}[chapter]
\theoremstyle{plain}
\newtheorem{lema}{Lemma}[chapter]
\theoremstyle{plain}
\newtheorem{cor}{Corollary}[chapter]
\theoremstyle{plain}
\newtheorem{resu}{Result}[chapter]
%
% Hyphenation for some words
%
\hyphenation{res-pec-tively}
\hyphenation{mono-ti-ca-lly}
\hyphenation{hypo-the-sis}
\hyphenation{para-me-ters}
\hyphenation{sol-va-bi-li-ty}
%
%
\begin{document}
% UoB guidlines:
%
% Preliminary pages
% 
% The five preliminary pages must be the Title Page, Abstract, Dedication
% and Acknowledgements, Author's Declaration and Table of Contents.
% These should be single-sided.
% 
% Table of contents, list of tables and illustrative material
% 
% The table of contents must list, with page numbers, all chapters,
 % sections and subsections, the list of references, bibliography, list of
% abbreviations and appendices. The list of tables and illustrations
% should follow the table of contents, listing with page numbers the
% tables, photographs, diagrams, etc., in the order in which they appear
% in the text.
% 
\frontmatter
\pagenumbering{roman}
%
%
% File: Title.tex
% Author: V?ctor Bre?a-Medina
% Description: Contains the title page
%
% UoB guidelines:
% 
% At the top of the title page, within the margins, the dissertation should give the title and, if 
% necessary, sub-title and volume number. If the dissertation is in a language other than English, the 
% title must be given in that language and in English. The full name of the author should be in the 
% centre of the page. At the bottom centre should be the words ?A dissertation submitted to the 
% University of Bristol in accordance with the requirements for award of the degree of ? in the 
% Faculty of ...?, with the name of the school and month and year of submission. The word count of 
% the dissertation (text only) should be entered at the bottom right-hand side of the page.
%
%
\begin{titlingpage}
\begin{SingleSpace}
\calccentering{\unitlength} 
\begin{adjustwidth*}{\unitlength}{-\unitlength}
\vspace*{13mm}
\begin{center}
\rule[0.5ex]{\linewidth}{2pt}\vspace*{-\baselineskip}\vspace*{3.2pt}
\rule[0.5ex]{\linewidth}{1pt}\\[\baselineskip]
{\HUGE What Does It Mean To Be Reactive?}\\[4mm]
{\Large \textit{Crushing opinions with the power of Mathematics}}\\
\rule[0.5ex]{\linewidth}{1pt}\vspace*{-\baselineskip}\vspace{3.2pt}
\rule[0.5ex]{\linewidth}{2pt}\\
\vspace{6.5mm}
{\large By}\\
\vspace{6.5mm}
{\large\textsc{Eddy Bertoluzzo}}\\
\vspace{11mm}
\includegraphics[scale=0.6]{logos/TU_Delft_logo_RGB}\\
\vspace{6mm}
{\large Faculty of Electrical Engineering, Mathematics and Computer Science\\
\textsc{Delft University of Technology}}\\
\vspace{11mm}
\begin{minipage}{10cm}
A dissertation submitted to Delft University of Technology in accordance with the requirements of the \textsc{Master of Science} degree in the Faculty of Electrical Engineering, Mathematics and Computer Science.
\end{minipage}\\
\vspace{9mm}
{\large\textsc{October 2016}}
\vspace{12mm}
\end{center}
\begin{flushright}
%{\small Word count: ten thousand and four}
\end{flushright}
\end{adjustwidth*}
\end{SingleSpace}
\end{titlingpage}
\clearemptydoublepage
%
%
% File: abstract.tex
% Author: V?ctor Bre?a-Medina
% Description: Contains the text for thesis abstract
%
% UoB guidelines:
%
% Each copy must include an abstract or summary of the dissertation in not
% more than 300 words, on one side of A4, which should be single-spaced in a
% font size in the range 10 to 12. If the dissertation is in a language other
% than English, an abstract in that language and an abstract in English must
% be included.

\chapter*{Abstract}
\begin{SingleSpace}
\initial{W}hat does it mean to be Reactive? The concept of Reactive Programming has gained much traction in the last few years as a paradigm well-suited for the development of asynchronous  event-driven applications. Unfortunately, Reactive Programming has been at the center of much discussion, if not confusion, with regards to its definition, properties and identifying principles. In this work we are going to wield the most powerful tool available to software engineers, mathematics, in order to formally derive the reactive types and bring clarity to this much opinionated topic.
\end{SingleSpace}
\clearpage
\clearemptydoublepage
%
%
% file: dedication.tex
% author: V?ctor Bre?a-Medina
% description: Contains the text for thesis dedication
%

\chapter*{Dedication and acknowledgements}
\begin{SingleSpace}
\initial{H}ere goes the dedication.
\end{SingleSpace}
\clearpage
\clearemptydoublepage
%
%%
% File: declaration.tex
% Author: V?ctor Bre?a-Medina
% Description: Contains the declaration page
%
% UoB guidelines:
%
% Author's declaration
%
% I declare that the work in this dissertation was carried out in accordance
% with the requirements of the University's Regulations and Code of Practice
% for Research Degree Programmes and that it has not been submitted for any
% other academic award. Except where indicated by specific reference in the
% text, the work is the candidate's own work. Work done in collaboration with,
% or with the assistance of, others, is indicated as such. Any views expressed
% in the dissertation are those of the author.
%
% SIGNED: .............................................................
% DATE:..........................
%
\chapter*{Author's declaration}
\begin{SingleSpace}
\begin{quote}
\initial{I} declare that the work in this dissertation was carried out in accordance with the requirements of  the University's Regulations and Code of Practice for Research Degree Programmes and that it  has not been submitted for any other academic award. Except where indicated by specific  reference in the text, the work is the candidate's own work. Work done in collaboration with, or with the assistance of, others, is indicated as such. Any views expressed in the dissertation are those of the author.

\vspace{1.5cm}
\noindent
\hspace{-0.75cm}\textsc{SIGNED: .................................................... DATE: ..........................................}
\end{quote}
\end{SingleSpace}
\clearpage
\clearemptydoublepage
%
\renewcommand{\contentsname}{Table of Contents}
\maxtocdepth{subsection}
\tableofcontents*
\addtocontents{toc}{\par\nobreak \mbox{}\hfill{\bf Page}\par\nobreak}
\clearemptydoublepage
%
%\listoftables
%\addtocontents{lot}{\par\nobreak\textbf{{\scshape Table} \hfill Page}\par\nobreak}
%\clearemptydoublepage
%
%\listoffigures
%\addtocontents{lof}{\par\nobreak\textbf{{\scshape Figure} \hfill Page}\par\nobreak}
%\clearemptydoublepage
%
%
% The bulk of the document is delegated to these chapter files in
% subdirectories.
\mainmatter
%
%\import{chapters/chapter01/}{chap01.tex}
\let\textcircled=\pgftextcircled
\chapter{Reactive Programming}
\label{chap:reactiveprogramming}

\epigraph{\hspace{4ex}\textit{The cold winds are rising in the North... Brace yourselves, winter is coming.}}{--- George R.R. Martin,\\ \textit{A Game of Thrones}}


In this chapter we are going to introduce the concept of reactive programming and motivate it's importance and relevance with regards to modern applications and the type of problems developers have to face nowadays. We are then going to introduce the most popular commercial libraries that claim to solve the reactive problem, with the purpose of giving the reader some context for our discussion and motivating the need for a mathematical formalization that abstracts over the class of problems these implementations set out to address.

\section{The Essence of Reactive Programs}

The use of the term reactive program in scientific literature is dated back to the mid-sixties\cite{scopus-reactive}. A relevant and insightful definition was given by G. Berry in 1991\cite{berry1991reactive} as he describes reactive programs in relation to their dual counterparts, interactive programs:

\begin{quote}
\hspace{4ex}``\textit{Interactive programs} interact at their own speed with users or with other programs; from a user point of view, a time-sharing system is interactive. \textit{Reactive programs} also maintain a continuous interaction with their environment, but at a speed which is determined by the environment, not by the program itself.''
\end{quote}

Interactive programs concretize the idea of a pull-based model of computation, where the program - the consumer in this case - has control over the speed at which data will be requested and handled. A perfect example of an interactive program is a control-flow structure such as a for-loop iterating over a collection of data: the program is in control of the speed at which data is retrieved from the containing collection and will request the next element only after it is done handling the current one.

Reactive programs, on the contrary, embody the idea of a push-based - or event-driven - model of computation, where the speed at which the program interacts with the environment is determined by the environment rather than the program itself. In other words, it is now the producer of the data - i.e. the environment - who determines the speed at which events will occur whilst the program's role reduces to that of a silent observer that will react upon receiving events. Standard example of such systems are GUI applications dealing with various events originating from user input - e.g. mouse clicks, keyboard button presses - and programs dealing with stock market, social media or any other kind of asynchronous updates.  

\section{Why Reactive Programming Matters}
\label{sec:whyrpmatters}

Considering the definition and examples of reactive programs we analyzed in the previous section, let's now try to formalize the class of problems the reactive programming paradigm is specifically well-suited for.

The table below provides a collection of types offered by common programming languages for handling data, parameterized over two variables: the size of the data, either one or multiple values, and the way data is handled, either by synchronous or asynchronous computations\cite{meijer2015spicing}.

\begin{center}
    \begin{tabular}{| l | l | l |}
    \hline
    & \textbf{One} & \textbf{Many} \\ 
    \hline
	\textbf{Sync} & \code{a} & \code{Iterable a} \\ 
	\hline
	\textbf{Async} & \code{Future a} & \textit{Reactive Programming} \\ 
	\hline
    \end{tabular}
    %\caption{The four fundamental effects}
\end{center}

The first row shows that synchronous functions come in two flavors: classic functions that return a single value of type \code{a} and functions that produce a collection of results of type \code{a}, abstracted through the \code{Iterable a} interface (See section \ref{sec:iterables}). These types of functions embody the standard imperative, pull-based approach to programming, where a call to a function/method synchronously blocks until a result is produced. 

Moving on to the second row, we encounter \code{Future a}, an interface representing an asynchronous computation that, at a certain point in the future, will result in a value of type \code{a}. Futures are generally created by supplying two callbacks together with the asynchronous computation, one to be executed in case of success and the other one in case of error. 

Programming languages, however, are not as well equipped when it comes to handling asynchronous computations resulting in multiple values - i.e. push-based collections. The issue lies in the fact that the program's control flow is dictated by the environment rather than the program itself - i.e. inversion of control -, making it very hard to model such problems with commonly known control structures, which are optimized for sequential models of computation. Traditional solutions typically involve developers manually trying to compose callbacks by explicitly writing CPS (continuation passing style) code\cite{meijer2015spicing}, resulting in what it's commonly referred to as \textit{Callback Hell}\cite{edwards2009coherent}.

The aforementioned class of problems reflects the definition of reactive programs we analyzed in the previous section, where the environment asynchronously - i.e. at its own speed - pushes multiple events to the program. The reactive programming paradigm sets out to provide interfaces and abstractions to facilitate the modeling of such problems as push-based collections. 

However, interfaces are only as good as the implementations that back them up. In the next section we are going to discuss and analyze various APIs and libraries that claim to embody the reactive paradigm, motivating our need for a mathematical formalization to aid in unifying these different approaches under a single set of interfaces, eliminating doubts and confusion when it comes to the essence of reactive programming. 

\section{Reactive Programming in the Real World}

Talk about FRP, Reactive Streams, Rx and other stuff. CIte survey paper
Claim that Observable is dual of Iterator in order to justify the why.

\subsection{Reactive Extensions}
\label{subsec:rx}

\subsection{Reactive Streams}
Try to define itself as reactive programming, mathematics prove the claim wrong, they implement asynchronous iterables. This does not make them useless, yet they are not a good solution for the problem we are trying to solve.

\todo[inline]{Probably put the derivation of Reactive Streams in the appendix, it's not really the focus, I guess.}

\subsection{Functional Reactive Programming}
\label{frp}
Much researched topic, beautiful in theory, yet based on continuous time, un-achievable with computers where time is inherently discrete. Connel Elliot, the inventor of FRP, says that commercial implementations are far from the original theory and do not really represent the essence of FRP
\clearemptydoublepage
%
%
% And the appendix goes here
\appendix
%\import{chapters/appendices/}{app0A.tex}
\chapter{Appendix A}
\label{app:a}

\section{Duality - Category Theory}

Duality is an element of paramount importance in research we presented. Up until this moment we have only dealt with an informal definition of duality but its roots go very deep in the field of mathematics and category theory in particular.

Given the formal definition of a category, composed of objects $$A,\; B,\; C,\; ...$$ and arrows $$f,\; g,\; h,\; ... $$ and operations $$ dom(f) \quad cod(f) \quad id(A) \quad g \circ f $$ 

Given any sentence $ \Sigma $, we can create its dual $ \Sigma^{op} $ by interchanging $$dom \to cod$$ and the order of composition $$ g \circ f \to f \circ g $$ 

It is clear how $ \Sigma^{op} $ is a well formed sentence in the language of category theory $ CT $.

The principle of - formal - duality then tells us that for any sentence $ \Sigma $ in the language of category theory, if $ \Sigma $ follows from the axioms defined for categories, then so does $ \Sigma^{op}$. 

$$ CT \vdash \Sigma \implies CT \vdash \Sigma^{op} $$

From a visual perspective, this definition boils down to reversing the order of composition of the arrows. 

\begin{center}
\begin{minipage}{.3\textwidth}
\begin{tikzpicture}
\large
  \matrix (m) [matrix of math nodes,row sep=3em,column sep=4em,minimum width=2em]
  { A & B \\
       & C \\};
  \path[-stealth]
    (m-1-1) edge node [above] {$f$} (m-1-2)
                 edge node [below] {$g \circ f\qquad$} (m-2-2)
    (m-1-2) edge node [right] {$g$} (m-2-2);
\end{tikzpicture}
\end{minipage}
\begin{minipage}{.3\textwidth}
\begin{tikzpicture}
\large
  \matrix (m) [matrix of math nodes,row sep=3em,column sep=4em,minimum width=2em]
  { A & B \\
       & C \\};
  \path[-stealth]
    (m-1-2) edge node [above] {$f$} (m-1-1)
    (m-2-2) edge node [below] {$f \circ g\qquad$} (m-1-1)
                 edge node [right] {$g$} (m-1-2);
\end{tikzpicture}
\end{minipage}
\end{center}

\section{Products \& Coproducts}

In category theory, the product of two objects in a category is the most general object which admits a morphism to each of the ones that compose it. The notion of product aims at capturing the essence of more specialized definitions found in various categories and areas of mathematics. The easiest way understand this construct is to start from the cartesian products in the category of sets.

Given sets $A$ and $B$, let us define the cartesian product as the set $$A \times B\; =\; \{\; (a, b)\; |\; a \in A,\; b \in B\; \} $$ There are two coordinate projections
\begin{center}
\begin{tikzpicture}
  \matrix (m) [matrix of math nodes,column sep=4em,minimum width=2em] { 
	  A & {A \times B} & B \\
	  & & \\
  };
  \path[-stealth]
    (m-1-2) edge node [above] {$fst$} (m-1-1)
                 edge node [above] {$snd$} (m-1-3);
\end{tikzpicture}
\end{center}
where $$ fst\; (a, b) = a \qquad snd\; (a, b) = b $$ 

It follows that given any element $ c \in A \times B $ $$ c\; =\; (\; fst\; (c),\; snd\; (c)\; )$$ 
The following diagram captures the essence of cartesian products.

\begin{center}
\begin{tikzpicture}
  \matrix (m) [matrix of math nodes,row sep=7em,column sep=6em,minimum width=2em] { 
	 {} & 1 & {} \\
	  A & {A \times B} & B \\
  };
  \path[-stealth]
  	(m-1-2) edge node [left] {$a\quad$} (m-2-1)
                 edge node [right] {$\quad b$} (m-2-3)
                 edge node [right] {$(a, b)$} (m-2-2)
    (m-2-2) edge node [below] {$fst$} (m-2-1)
                 edge node [below] {$snd$} (m-2-3);
\end{tikzpicture}
\end{center}

The definition of categorical products can be derived from generalizing the elements in the previous definition.

Coproducts are the dual notion of categorical products, representing the least general object to which the objects in the family admit a morphism. Within the context of set theory, the represent the disjoint union of sets.

Given sets $A$ and $B$, let us define their disjoint union as the set $$A + B\; =\; \{\; (a, 1)\; |\; a \in A\} \cup (b, 2)\; |\; b \in B\} $$ There are two injection functions
\begin{center}
\begin{tikzpicture}
  \matrix (m) [matrix of math nodes,column sep=4em,minimum width=2em] { 
	  A & {A + B} & B \\
	  & & \\
  };
  \path[-stealth]
    (m-1-1) edge node [above] {$left$} (m-1-2)
    (m-1-3) edge node [above] {$right$} (m-1-2);
\end{tikzpicture}
\end{center}
where $$ left\; (a) = (a,\; 1) \qquad right\; (b) = (b,\; 2) $$ 

The essence is captured by the following diagram, where we have simply reversed the arrows.

\begin{center}
\begin{tikzpicture}
  \matrix (m) [matrix of math nodes,row sep=7em,column sep=6em,minimum width=2em] { 
	 {} & 1 & {} \\
	  A & {A + B} & B \\
  };
  \path[-stealth]
  (m-2-1) edge node [left] {$f\quad$} (m-1-2)
  (m-2-3) edge node [right] {$\quad g$} (m-1-2)
  (m-2-2) edge node [right] {$[f,g]$} (m-1-2)
  (m-2-1) edge node [below] {$left$} (m-2-2)
  (m-2-3) edge node [below] {$right$} (m-2-2);
\end{tikzpicture}
\end{center}

where $$ [f, g](x,\delta) = \begin{cases} f(x), & \delta = 1 \\ g(x), & \delta = 2 \end{cases} $$

Withing the context of our research and the Haskell programming language, tuples represent products and the Either monad represents coproducts.

\section{Functors}

In category theory a functor is a mapping from one category to another or a homomorphism of categories where certain laws hold. 

Functors span thorough a large number of categories and are more and more common in modern programming languages. In Haskell the Functor class and its laws are defined as follows.

\hminted
class Functor f where
    fmap :: (a -> b) -> f a -> f b
  
fmap id      = id
fmap (p . q) = (fmap p) . (fmap q)
\end{minted}

\section{Covariance \& Contravariance}

\section{Continuations}

\section{Curry \& Uncurry}

Currying is the technique of transforming a function taking multiple arguments as input in one taking only the first argument and returning a function that takes the remainder of the arguments and returns the result of the initial function. Uncurrying is the opposite of currying, taking a curried function into one that accepts multiple arguments as input.

\hminted
curry   :: ((a, b) -> c) -> (a -> b -> c)
uncurry :: (a -> b -> c) -> (a, b) -> c

f :: a -> b -> c
g :: (a, b) -> c

f = curry g
g = uncurry f

-- f x y = g (x, y)
-- curry . uncurry = id
\end{minted}










\
\clearemptydoublepage
%
% Apparently the guidelines don't say anything about citations or
% bibliography styles so I guess we can use anything.
\backmatter
\bibliographystyle{siam}
\refstepcounter{chapter}
\bibliography{thesisbiblio}
\clearemptydoublepage
%
% Add index
%\printindex
%   
\end{document}